\documentclass[a4paper]{article}
\usepackage{fullpage}
\usepackage[margin=63pt]{geometry}
\usepackage{graphics}
\usepackage{xcolor}
\usepackage{graphicx}
\usepackage{subcaption}
\usepackage{amssymb}
\usepackage{pdfpages}
\usepackage{booktabs}

\input{../stachnisslab-math.tex}
%% Standard latex definitions used at the Stachniss-Lab

\usepackage{graphics}           
\usepackage{times}              
\usepackage{amsmath}            
\usepackage{amssymb}            
\usepackage{graphicx}
\usepackage{algorithm}
\usepackage[noend]{algpseudocode}
\usepackage{booktabs}
\usepackage{balance}
\usepackage{color}
\usepackage[nocompress]{cite} % sorts citations automatically, e.g., "[1],[2],[3]" instead of "[2],[3],[1]".
\definecolor{instructioncolor}{rgb}{.5,.5,.5}
\usepackage{multirow} 
\usepackage{url} 
\newcommand{\instruction}[1]{\textcolor{instructioncolor}{#1}}

% make caption font small for better separation of figures and text
\usepackage[font=small]{caption}
\usepackage{svg}
 %% Key definitions for text elements. USE ONLY THEM! Do not use naked \ref{}.
\def\secref#1{Sec.~\ref{#1}}
\def\figref#1{Fig.~\ref{#1}}
\def\tabref#1{Tab.~\ref{#1}}
\def\eqref#1{Eq.~(\ref{#1})}
\def\algref#1{Alg.~\ref{#1}}

%% Other useful macros
\newcommand\todo[1]{\textbf{[TODO: #1}]}

\makeatletter
\usepackage{xspace}
%% the \onedot macro is producing only one dot at line ends.
%% thus \etal. will not produce et al..
\DeclareRobustCommand\onedot{\futurelet\@let@token\@onedot}
\def\@onedot{\ifx\@let@token.\else.\null\fi\xspace}
\def\eg{e.g\onedot} \def\Eg{E.g\onedot}
\def\ie{i.e\onedot} \def\Ie{I.e\onedot}
\def\cf{cf\onedot} \def\Cf{Cf\onedot}
\def\etc{etc\onedot} \def\vs{vs\onedot}
\def\wrt{w.r.t\onedot} \def\dof{d.o.f\onedot}
%% Cyrill does not like emph...
%\def\etal{\emph{et al}\onedot}
\def\etal{{et al}\onedot}
\makeatother

\def\etalcite#1{\etal~\cite{#1}}

\usepackage{array}
%% this allows to use something like p{2cm} as column type, but with left, center, and right alignment
\newcolumntype{L}[1]{>{\raggedright\let\newline\\\arraybackslash\hspace{0pt}}m{#1}}
\newcolumntype{C}[1]{>{\centering\let\newline\\\arraybackslash\hspace{0pt}}m{#1}}
\newcolumntype{R}[1]{>{\raggedleft\let\newline\\\arraybackslash\hspace{0pt}}m{#1}}




\def\etalcite#1{\etal~\cite{#1}}
\newcommand\circledmark[1][red]{%
  \ooalign{%
    \hidewidth
    \kern0.65ex\raisebox{-0.9ex}{\scalebox{3}{\textcolor{#1}{\textbullet}}}
    \hidewidth\cr
    $\checkmark$\cr
  }%
}

\definecolor{brewerGreen0}{HTML}{E5F5F9}
\definecolor{brewerGreen1}{HTML}{99D8C9}
\definecolor{brewerGreen2}{HTML}{2CA25F}
\definecolor{brewerCyan0}{HTML}{ECE2F0}
\definecolor{brewerCyan1}{HTML}{A6BDDB}
\definecolor{brewerCyan2}{HTML}{1C9099}
\definecolor{brewerCyan9}{HTML}{023858}
\definecolor{brewerGrey0}{HTML}{F0F0F0}
\definecolor{brewerGrey1}{HTML}{BDBDBD}
\definecolor{brewerGrey2}{HTML}{636363}

\definecolor{commentColor}{HTML}{000000}
\definecolor{revisionColor}{HTML}{0238A8}
\definecolor{answerBoxColor}{HTML}{006D2C}

\usepackage[colorlinks=true,linkcolor=answerBoxColor,citecolor=answerBoxColor,urlcolor=answerBoxColor,pagebackref]{hyperref}

\usepackage{tcolorbox}
\def\answer#1#2{\begin{tcolorbox}[width=\textwidth,colback={white},title={\textbf{Answer
 #1}},colbacktitle=answerBoxColor,coltitle=white] #2 \end{tcolorbox} }

\newcounter{todonum}
\newcommand\resettodonum{\setcounter{todonum}{0}}
\newcommand\mytodo{%
  \stepcounter{todonum}%
  \textcolor{red}{\;TODO \thetodonum}
}

\newcommand\boldblue[1]{\textcolor{blue}{\textbf{#1}}}
\newcommand\boldred[1]{\textcolor{red}{\textbf{#1}}}
\newcommand{\reviewed}[1]{\textcolor{blue}{#1}}

\title{
  \textbf{Anonymous Manuscript Revision}
  \\[\baselineskip] Authors' Response to AE and Reviewers
}
\author{Anonymous Authors}

\usepackage{titling}
\setlength{\droptitle}{-0.5in}

\begin{document}
\maketitle
\resettodonum

We thank the Associate Editor and both reviewers for their careful and constructive feedback.
In the revised manuscript, we addressed all requested points and added new experiments for fairness,
robustness, and runtime reporting.
For readability, modifications in the revised manuscript are highlighted in \reviewed{blue}.

\paragraph{Summary of major revisions}
\begin{itemize}
  \item \textbf{Fairer benchmarking:} added \textbf{AMCL+NoisyGPS} (Kalman fusion) baseline in both traverses and reported multiseed results (seeds 11, 22, 33).
  \item \textbf{Clearer metrics:} report raw APE and \textbf{RPE at 2 m / 5 m / 10 m} explicitly, plus cross-track error, row-correct fraction, and row-switch events.
  \item \textbf{Semantic-wall clarity:} expanded the methodology and equations for ray casting and semantic-wall likelihood terms, and added a point-only non-wall formulation for comparison.
  \item \textbf{Robustness to semantic/map degradation:} added two studies: (A) random semantic detection drop (20\%, 40\%) and (B) section landmark removal (30\%, 50\%) with recovery metrics.
  \item \textbf{Runtime evaluation:} added repeated runtime profiling with full-pipeline throughput and component-level timings.
  \item \textbf{Experimental reproducibility details:} expanded hardware/sensor/setup details (Thorvald platform, Intel RealSense D435i, SICK Tim7 2D LiDAR, RTK-GNSS), map usage, and fixed parameters.
  \item \textbf{Figure readability:} regenerated trajectory figures with larger labels/legends, clearer line rendering, landmark/semantic-wall overlays, and start/end markers.
  \item \textbf{Text fixes:} corrected all minor typographical/grammar issues raised by the reviewers.
\end{itemize}

\section*{Associate Editor}
\paragraph{Report}
{\color{commentColor}
(1) Experimental evaluation is difficult to interpret; fairness concern because the proposed method uses GPS while some baselines do not; clarify RPE as RPE@X meters.\\
(2) The advantages of semantics (especially semantic walls) are not sufficiently clear methodologically and quantitatively.\\
(3) Plot readability needs improvement (labels/lines difficult to see in key figures).
}

\vspace{8pt}
\answer{to Associate Editor}{
Thank you for the detailed guidance. We revised the paper along exactly these three axes:
\begin{itemize}
  \item \textbf{Interpretability and fairness:} we added a fused \emph{AMCL+NoisyGPS} baseline and reran experiments on both traverses with three seeds. We now report raw APE and \emph{RPE@2m/@5m/@10m}. For example, on run2, SLPF obtains raw APE $=1.24\pm0.04$ m, versus $3.38\pm0.91$ m for AMCL+NoisyGPS and $3.50\pm0.94$ m for AMCL.
  \item \textbf{Semantic contribution clarity:} we expanded the semantic-wall formulation and added explicit ablations. In the main ablation, removing semantic information degrades row consistency and geometry (row-correct from $0.732$ to $0.673$, cross-track mean from $1.263$ m to $1.499$ m).
  \item \textbf{Plot clarity:} we regenerated the affected figures with larger fonts/legends, thicker trajectories, consistent scales, and overlays of landmarks/semantic walls to make row-switch behaviour easier to interpret.
\end{itemize}
}

\newpage
\section*{Reviewer 1}

\paragraph{Comment 1.1}
{\color{commentColor}
The introduction/related work underplays prior fusion methods (LiDAR+vision/semantics). Please expand prior work discussion, clarify novelty and assumptions, and justify the 2D LiDAR choice versus 3D LiDAR.
}
\answer{1.1}{
We expanded the introduction and related-work discussion to better position prior fusion/localization methods and to state our novelty more explicitly.
We now clarify that our main technical contribution is not generic fusion, but a semantic-particle-filter design that uses semantic wall constraints plus adaptive GNSS weighting to handle row aliasing in repetitive vineyards.
We also added a dedicated discussion on the 2D-vs-3D LiDAR design tradeoff: expected gains with 3D sensing, what would change in our pipeline, and why the current 2D setup was chosen for this platform and deployment context.
}

\paragraph{Comment 1.2}
{\color{commentColor}
Abstract/Introduction clarity: explicitly mention camera-based semantics (RGB-D), define AMCL at first use, provide quantitative effect sizes, state LiDAR type, and specify required navigation accuracy for vineyard autonomy.
}
\answer{1.2}{
We revised the abstract and introduction accordingly:
\begin{itemize}
  \item semantics are explicitly described as camera/RGB-D driven;
  \item AMCL is expanded on first use;
  \item quantitative headline results are included;
  \item LiDAR modality is stated in the abstract;
  \item the introduction now states the operational row-following requirements and links them to the reported metrics (raw APE/RPE/cross-track/row-correct).
\end{itemize}
}

\paragraph{Comment 1.3}
{\color{commentColor}
Method reproducibility: clarify data flow/parameters, calibration and time synchronization, and map construction assumptions (cost, maintenance, partial/missing landmarks).
}
\answer{1.3}{
We expanded the methodology and setup descriptions for reproducibility:
\begin{itemize}
  \item clearer end-to-end data flow from RGB-D instance masks to BEV semantic labels and PF likelihood updates;
  \item explicit fixed parameterization (including GNSS/semantic/corridor/background terms and particle settings);
  \item explicit statement that synchronized RGB-D/LiDAR/RTK streams are used, with known camera--robot extrinsics in projection;
  \item added map construction discussion, including behaviour under partial/perturbed maps and maintenance implications.
\end{itemize}
We also added dedicated robustness experiments under map degradation (see Comment 1.5/2.1 responses).
}

\paragraph{Comment 1.4}
{\color{commentColor}
Experimental setup and baselines: report vineyard size/path length and LiDAR details; clarify RTAB-Map settings; consider a modern visual(-inertial) SLAM baseline; improve qualitative trajectory interpretation with landmark/wall overlays and row-switch focus.
}
\answer{1.4}{
We revised the experimental setup section with explicit site and sensor details (10 rows, mean row length, area/footprint, platform and sensor list including SICK Tim7 2D LiDAR).
We also clarified baseline protocol/settings:
\begin{itemize}
  \item RTAB-Map RGB/RGBD run with default configuration files;
  \item AMCL configured with the same 5 m sensing range used by SLPF for comparability;
  \item added AMCL+NoisyGPS fusion baseline for fairness;
  \item multiseed evaluation on two traverses with aggregate statistics.
\end{itemize}
For qualitative clarity, we regenerated trajectory figures with semantic-wall/landmark overlays, clearer labels and line styles, and start/end markers.
In addition, we executed ORB-SLAM3 runs in our evaluation pipeline and report them as an additional reference baseline in the revision package.
}

\paragraph{Comment 1.5}
{\color{commentColor}
Failure modes should be quantified, e.g., degrade detections or mask map portions and report performance deltas.
}
\answer{1.5}{
We added two new robustness studies on run1 (3 seeds):
\begin{itemize}
  \item \textbf{Detection Dropout}: random semantic detection removal at 20\% and 40\%.
  \item \textbf{Section Landmark Removal}: remove 30\% and 50\% landmarks in one contiguous map band and evaluate recovery after leaving that region.
\end{itemize}
Results show graceful degradation and recovery. For example, baseline raw APE is $1.02\pm0.03$ m; with 40\% detection dropout it is $1.08\pm0.11$ m. Under 50\% section removal, post-section recovery remains successful in all runs (recovery rate $=100\%$) with mean recovery distance $11.79\pm4.46$ m.
These additions directly quantify failure sensitivity and recovery behaviour.
}

\newpage
\section*{Reviewer 2}

\paragraph{Comment 2.1}
{\color{commentColor}
The semantic-wall contribution is not sufficiently clear in text/equations; clarify how virtual walls are fused with real observations and compare to using only real landmarks.
}
\answer{2.1}{
We substantially revised the method section to clarify this point.
The revised manuscript now details:
\begin{itemize}
  \item semantic-wall map construction from adjacent landmarks;
  \item ray-casting against wall segments to obtain class/range predictions;
  \item class-aware semantic/background likelihood terms and their fusion with GNSS/corridor terms.
\end{itemize}
To address the quantitative part, we added an explicit non-wall point-based formulation and ablation.
Compared with the full model, removing semantics worsens row-level robustness (row-correct $0.732\rightarrow0.673$ and cross-track $1.263\rightarrow1.499$ m in the main ablation).
This clarifies both \emph{how} walls are fused and \emph{what} they contribute empirically.
}

\paragraph{Comment 2.2}
{\color{commentColor}
For fairness, baselines such as AMCL/RTAB-Map should also include GPS observations since the proposed method uses GPS.
}
\answer{2.2}{
Agreed. We added a fused \textbf{AMCL+NoisyGPS} baseline (constant-velocity Kalman fusion of AMCL and noisy GNSS positions) and evaluated it in the same multiseed protocol.
This directly addresses the fairness concern.
Representative results:
\begin{itemize}
  \item run1: SLPF raw APE $=1.14\pm0.07$ m vs AMCL+NoisyGPS $1.89\pm0.04$ m;
  \item run2: SLPF raw APE $=1.24\pm0.04$ m vs AMCL+NoisyGPS $3.38\pm0.91$ m.
\end{itemize}
We now report all methods with raw APE, RPE@2/5/10m, cross-track, row-correct fraction, and row-switch events.
}

\paragraph{Comment 2.3}
{\color{commentColor}
Please include runtime evaluation to show real-time capability.
}
\answer{2.3}{
We added a dedicated runtime profiling experiment (3 trials, warmup excluded).
On the reported hardware (Threadripper 2950X + RTX 4060 Ti/3070), we obtain:
\begin{itemize}
  \item processed-frame throughput: $13.04\pm0.21$ Hz,
  \item effective input-rate throughput with stride 4: $52.18\pm0.84$ Hz.
\end{itemize}
We also report component-level timing breakdown (semantic inference, measurement update, LiDAR association, motion update, resampling), enabling reproducibility and deployment-oriented interpretation.
}

\paragraph{Comment 2.4}
{\color{commentColor}
Minor points: increase axis/legend readability; fix Figure 3 caption punctuation; correct grammar/typographical issues in Sections III.C, III.D.2, and IV.B.
}
\answer{2.4}{
All minor issues were corrected in the revision:
\begin{itemize}
  \item increased font sizes and improved readability in the affected figures;
  \item corrected the Figure 3 caption punctuation;
  \item fixed the grammatical/typographical errors identified in Sections III.C, III.D.2, and IV.B.
\end{itemize}
}

\end{document}
